% SOURCE_STATEMENT: "Can Herbst's technical regularity assumption in the extension of Haag–Ruelle scattering theory to theories with massless particles be replaced by a physically motivated criterion?"
% Preparation and final self-audit: CONTRIBUTING.md prompts C and D; source/audit artifacts retained in the local physics_statements_2026_10_11 batch.
% !TeX program = lualatex
\documentclass[11pt]{article}
\usepackage[margin=25mm]{geometry}
\usepackage{fontspec}
\setmainfont{Latin Modern Roman}
\usepackage{amsmath,amssymb,mathtools}
\usepackage{unicode-math}
\setmathfont{Latin Modern Math}
\usepackage{xurl,hyperref}
\hypersetup{colorlinks=true,urlcolor=blue}
\setcounter{secnumdepth}{0}
\setlength{\parindent}{0pt}
\setlength{\parskip}{5pt}
\setlength{\emergencystretch}{3em}
\newcommand{\sref}[2]{\hyperlink{source:#1:#2}{[S#2]}}
\newcommand{\cataloguescope}[1]{\paragraph{Recorded target.}#1}
\newcommand{\R}{\mathbb R}
\newcommand{\C}{\mathbb C}
\newcommand{\N}{\mathbb N}
\newcommand{\Z}{\mathbb Z}
\newcommand{\Q}{\mathbb Q}
\newcommand{\F}{\mathbb F}
\newcommand{\eps}{\varepsilon}
\begin{document}
\section{Physical Regularity Criteria for Massless Scattering Theory}\label{statement}

\subsection{Definitions and mathematical statement}

\paragraph{Shared notation.}$\mathbb N=\{1,2,\ldots\}$; the standard number fields have their usual meanings. All additional conventions are stated locally.


\paragraph{Definitions and conventions.} Haag--Ruelle scattering theory constructs asymptotic multiparticle states from local observables in a relativistic quantum field theory with a vacuum and translation spectrum. A sharp-mass particle has a one-particle mass shell even when other excitations accumulate near it. Herbst's regularity assumption controls the unwanted spectral component near this mass shell sufficiently strongly for the scattering limits. The report concerns a theory that also has massless particles, while the scattered massive particles retain sharp masses.

\paragraph{Source scope and qualifications.} The report excludes electrically charged infraparticles through the sharp-mass hypothesis. The exact quantitative Herbst assumption is in the cited scattering article rather than in this short report passage; no invented exponent or spectral bound is substituted here.

\paragraph{Catalogue mathematical target.}
\paragraph{Catalogue wording retained for traceability.}
Can Herbst's technical regularity assumption in the extension of Haag–Ruelle scattering theory to theories with massless particles be replaced by a physically motivated criterion?

\paragraph{Saved literature context.} The supplied catalogue's assessment dated 2026-08-21 records: The report explicitly identifies replacing Herbst's technical assumption by a physically motivated criterion as an open problem. The assessment is retained as dated provenance; it is not a new comprehensive literature or open-status certification.
\subsection{Short English statement}
Replace the technical spectral regularity assumption in massless-background Haag--Ruelle scattering by a physically meaningful condition.
\subsection{Sources}
\begin{itemize}
\item[\hypertarget{source:statement:1}{S1}] ULAM.AI and the UnsolvedMath Contributors, UnsolvedMath, catalogue entry. Catalogue curation under CC BY 4.0; definitions and English exposition added locally.
\url{https://huggingface.co/datasets/ulamai/UnsolvedMath}.
\item[\hypertarget{source:statement:2}{S2}] Analysis and Quantum Theory (2005), Oberwolfach Reports, DOI 10.4171/owr/2005/44, pp. 2493--2518; official report source and mathematical context.
\url{https://doi.org/10.4171/owr/2005/44}.
\item[\hypertarget{source:statement:3}{S3}] Official EMS Press PDF of the same Oberwolfach report; PDF pages 7--8, printed pages 2499--2500.
\url{https://ems.press/content/serial-article-files/46017}.
\end{itemize}
\label{end:statement}
\end{document}
